\chapter{Master Problem Bank: Electrolysis, Faraday's Laws \& Overvoltage}
\label{chap:qb_electrolysis}

\begin{tcolorbox}[enhanced,colback=subtlegold,colframe=accentamber,arc=2mm,boxrule=1pt,
    title=\textbf{\large Part VI Organization: Electrolysis, Faraday's Laws \& Kinetics}]
This part compiles all questions and quantitative problems on preferential discharge, overpotential ($\eta$) on various electrode materials, Faraday's 1st and 2nd laws, current efficiency, and industrial electrolytic cells from all 8 books. Identical and redundant variants from other books are co-located beneath each primary problem with full source citations.
\end{tcolorbox}

\section{Subtopic 6.1: Preferential Discharge Theory \& Overpotential ($\eta$)}

\begin{problembox}
\textbf{Q.6.1} \hfill \textbf{[Primary Source: Neeraj Kumar Part 2 Ex-II Sec B Q4 / GRB  12.3]}
\label{q:qb-c6-01}

Predict the primary product liberated at the cathode and at the anode during the electrolysis of an aqueous $1\,\text{M}$ solution of \ce{NaCl} under the following two distinct experimental conditions:
\begin{enumerate}[label=(\alph*)]
    \item Using inert platinum (Pt) electrodes at low current density
    \item Using a liquid mercury (Hg) cathode and graphite anode (Castner--Kellner cell)
\end{enumerate}
Explain the physical role of hydrogen overvoltage on mercury ($\eta_{\ce{H2/Hg}} \approx 1.5\,\text{V}$) in determining the cathodic product.

\tcbline
\textbf{\color{accentamber}Cross-Book Redundant / Identical Variants:}
\begin{itemize}[leftmargin=*,itemsep=2pt]
    \item \textbf{[Variant v1: Pearson Ex-II Sec A Q56, p. 21]} During the electrolysis of aqueous \ce{NaCl} with mercury cathode, sodium is discharged instead of hydrogen because: (a) $E^\circ(\ce{Na+/Na}) > E^\circ(\ce{H+/H2})$ (b) hydrogen overvoltage on mercury is exceptionally high (c) sodium forms amalgam with mercury lowering its discharge potential (d) both (b) and (c).
    \item \textbf{[Variant v2: Narendra Avasthi Level-1 Q142, p. 20]} During electrolysis of aqueous \ce{CuSO4} solution using copper electrodes, the process taking place at the anode is: (a) evolution of \ce{O2} (b) evolution of \ce{SO2} (c) dissolution of copper into $\ce{Cu^2+}$ (d) deposition of sulfate.
    \item \textbf{[Variant v3: Cengage Theory p. 78]} Why does chlorine gas evolve at the anode during electrolysis of concentrated aqueous \ce{NaCl} even though thermodynamic standard reduction potential of oxygen ($E^\circ = 1.23\,\text{V}$) is lower than that of chlorine ($E^\circ = 1.36\,\text{V}$)?
\end{itemize}
\end{problembox}

\section{Subtopic 6.2: Faraday's 1st and 2nd Laws \& Current Efficiency}

\begin{problembox}
\textbf{Q.6.2} \hfill \textbf{[Primary Source: GRB Practice Problem 64, p. 52 / NA Level-1 Q58]}
\label{q:qb-c6-02}

A constant current of $3.0\,\text{A}$ is passed for $2.0\,\text{hours}$ through three electrolytic cells connected in series containing solutions of:
(A) \ce{AgNO3(aq)}, (B) \ce{CuSO4(aq)}, (C) \ce{CrCl3(aq)}.
Assuming $100\%$ current efficiency, calculate the mass deposited at the cathode in each cell:
\begin{enumerate}[label=(\alph*)]
    \item Mass of silver (Atomic mass $= 108.0$)
    \item Mass of copper (Atomic mass $= 63.5$)
    \item Mass of chromium (Atomic mass $= 52.0$)
\end{enumerate}

\tcbline
\textbf{\color{accentamber}Cross-Book Redundant / Identical Variants:}
\begin{itemize}[leftmargin=*,itemsep=2pt]
    \item \textbf{[Variant v1: Neeraj Kumar Part 2 Ex-I Q68, p. 38]} The same quantity of electrical charge that deposited $1.08\,\text{g}$ of silver from \ce{AgNO3} will deposit what mass of aluminium from molten \ce{AlCl3}? (a) $0.09\,\text{g}$ (b) $0.27\,\text{g}$ (c) $0.81\,\text{g}$ (d) $1.08\,\text{g}$.
    \item \textbf{[Variant v2: Essential  24.4 Solved Ex 3, p. 863]} What current strength in amperes is required to deposit $5.00\,\text{g}$ of gold from a solution of \ce{HAuCl4} in $1.0\,\text{hour}$? ($\text{At. mass of Au} = 197.0$, $n = 3$).
    \item \textbf{[Variant v3: Pearson Ex-I Q82, p. 7]} In an electroplating process, a current of $10.0\,\text{A}$ deposits $12.0\,\text{g}$ of zinc in $1.0\,\text{hour}$. The current efficiency ($\eta_{\text{curr}}$) is: (a) $98.5\%$ (b) $82.4\%$ (c) $65.2\%$ (d) $50.0\%$.
\end{itemize}
\end{problembox}

\begin{problembox}
\textbf{Q.6.3} \hfill \textbf{[Primary Source: Narendra Avasthi Level-2 Q10, p. 25 / PRSN Ex-II Q8]}
\label{q:qb-c6-03}

An aqueous solution of $0.1\,\text{M } \ce{CuSO4}$ ($1.0\,\text{L}$) is electrolyzed between inert Pt electrodes with a steady current of $9.65\,\text{A}$ for $1000\,\text{seconds}$.
\begin{enumerate}[label=(\alph*)]
    \item Calculate the mass of copper deposited at the cathode.
    \item Calculate the total volume of oxygen gas evolved at the anode measured at STP ($0^\circ\text{C}, 1\,\text{atm}$).
    \item Calculate the final $\text{pH}$ of the solution after electrolysis (assume volume remains constant at $1.0\,\text{L}$).
\end{enumerate}

\tcbline
\textbf{\color{accentamber}Cross-Book Redundant / Identical Variants:}
\begin{itemize}[leftmargin=*,itemsep=2pt]
    \item \textbf{[Variant v1: Cengage DPP 3.2 Q19, p. 4]} A current strength of $96.5\,\text{A}$ is passed for $10\,\text{seconds}$ through $1\,\text{L}$ of $0.1\,\text{M}$ aqueous \ce{CuSO4}. The pH of the resulting solution is: (a) $1.0$ (b) $2.0$ (c) $3.0$ (d) $7.0$.
    \item \textbf{[Variant v2: Neeraj Kumar Part 2 Ex-II Sec A Q34, p. 44]} When $100\,\text{mL}$ of $0.2\,\text{M } \ce{AgNO3}$ is electrolyzed until half the silver is deposited, how many moles of $\ce{H+}$ ions are produced in the electrolyte?
    \item \textbf{[Variant v3: GRB Practice Problem 68, p. 52]} Electrolysis of dilute \ce{H2SO4} for $1.0\,\text{hour}$ produces $112\,\text{mL}$ of \ce{O2} at STP at the anode. What volume of \ce{H2} is evolved at the cathode at the same temperature and pressure?
\end{itemize}
\end{problembox}
